% \makenoidxglossaries

%...UNIT 1: CONSTANT MOTION

\newglossaryentry{magnitude}{
    name={magnitude},
    description={size or amount}
}

\newglossaryentry{scalar}{
    name=scalar,
    description={a quantity that has magnitude but no direction}
}

\newglossaryentry{vector}{
    name={vector},
    description={a quantity that has both magnitude and direction}
}

\newglossaryentry{component}{
    name={component},
    description={independent horizontal ($x$) and vertical ($y$) projections of a two-dimensional vector}
}

\newglossaryentry{tail}{
    name={tail},
    description={the starting point of a vector}
}

\newglossaryentry{head}{
    name={head},
    description={the end point of a vector}
}


\newglossaryentry{head-to-tail method}{
    name={head-to-tail method},
    description={a method of adding vectors in which the tail of each vector is placed at the head of the previous vector}
}

\newglossaryentry{vector sum}{
    name={vector sum},
    description={the result of adding two or more vectors together, taking into account both their magnitudes and directions}
}

\newglossaryentry{position}{
    name={position},
    description={the location of an object at a particular time.}
}

\newglossaryentry{motion}{
    name={motion},
    description={a change in an object's position over time}
}

\newglossaryentry{reference frame}{
    name={reference frame},
    description={a coordinate system used to describe an object's position and motion}
}

\newglossaryentry{one-dimensional motion}{
    name={one-dimensional motion},
    description={motion in which an object moves along a straight line in a single direction}
}


\newglossaryentry{displacement}{
    name={displacement},
    description={the change in an object's position from its starting point to its ending point.}
}

\newglossaryentry{distance}{
    name={distance},
    description={the length of the path actually traveled between an initial and a final position}
}

\newglossaryentry{position vs. time graph}{
    name={position vs. time graph},
    description={a graph in which position is plotted on the vertical axis and time is plotted on the horizontal axis}
}

\newglossaryentry{speed}{
    name={speed},
    description={rate at which an object changes its location}
}

\newglossaryentry{average speed}{
    name={average speed},
    description={distance traveled divided by the time during which the motion occurs}
}

\newglossaryentry{velocity1}{
    name={velocity},
    description={the speed and direction of motion}
}

\newglossaryentry{velocity2}{
    name={velocity},
    description={the rate at which position changes with time}
}

\newglossaryentry{average velocity}{
    name={average velocity},
    description={displacement divided by the time during which the displacement occurs}
}

\newglossaryentry{velocity vs. time graph}{
    name={velocity vs.\,time graph},
    description={a graph in which velocity is plotted on the vertical axis and time is plotted on the horizontal axis}
}


\newglossaryentry{relative speed}{
    name={relative speed},
    description={how fast or slow an object appears to be moving to another object}
}

\newglossaryentry{relative velocity}{
    name={relative velocity},
    description={the rate at which an object changes position relative to another object}
}

%...UNIT 2: ACCELERATION

\newglossaryentry{acceleration}{
    name={acceleration},
    description={the rate at which velocity changes over time}
}

\newglossaryentry{acceleration2}{
    name={acceleration},
    description={the speeding up, slowing down, or changing direction of an object}
}

\newglossaryentry{average acceleration}{
    name={average acceleration},
    description={the change in velocity divided by the time interval over which it changed}
}

\newglossaryentry{acceleration vs. time graph}{
    name={acceleration vs.\,time graph},
    description={a graph in which acceleration is plotted on the vertical axis and time is plotted on the horizontal axis}
}

%...UNIT 3: FORCES

\newglossaryentry{system}{
    name={system},
    description={a defined object or group of objects chosen to be analyzed}
}

\newglossaryentry{force}{
    name=force,
    description={a push or pull acting on an object}
}

\newglossaryentry{net force}{
    name={net force},
    description={the vector sum of all forces acting on an object}
}

\newglossaryentry{external force}{
    name=external force,
    description={a force that originates outside of the object or system}
}

\newglossaryentry{net external force}{
    name={net external force},
    description={the vector sum of all external forces acting on an object}
}

\newglossaryentry{tension}{
    name=tension,
    description={a pulling force transmitted through a string, rope, or cable}
}

\newglossaryentry{friction}{
    name={frictional force},
    description={a force that opposes the relative motion between two surfaces in contact}
}

\newglossaryentry{normal force}{
    name={normal force},
    description={a force exerted by a surface on an object, perpendicular to the surface}
}

\newglossaryentry{gravity}{
    name={gravity},
    description={a force of attraction between objects due to their masses}
}

\newglossaryentry{applied force}{
    name={applied force},
    description={a force exerted on an object by contact with another object or person}
}

\newglossaryentry{gravitational force}{
    name={gravitational force}, %...MY DEFINITION
    description={the downward force on an object due to the attraction by the Earth %or other massive body
    }
}

\newglossaryentry{spring force}{
    name=spring force, %...From district slides
    description={a force applied from a spring when it is either compressed or stretched}
}

\newglossaryentry{free body diagram}{
    name=free body diagram,
    description={a diagram that represents an object as a single point and shows all external forces acting on it as arrows (vectors)}
}


\newglossaryentry{mass}{
    name=mass,
    description={the quantity of matter in an object}
}

\newglossaryentry{Newton's first law of motion}{
    name={Newton's first law of motion},
    description={an object at rest stays at rest, and an object in motion stays in motion at constant velocity, unless acted upon by a net external force}
}

\newglossaryentry{inertia}{
    name=inertia,
    description={the tendency of an object to resist any change in its state of motion}
}

\newglossaryentry{Newton's second law of motion}{
    name={Newton's second law of motion},
    description={the acceleration of an object is directly proportional to the net force acting on it and inversely proportional to its mass}
}

\newglossaryentry{Newton's third law of motion}{
    name={Newton's third law of motion},
    description={for every action force, there is an equal and opposite reaction force}
}

%...UNIT 4: ONE-DIMENSIONAL KINEMATICS

\newglossaryentry{kinematic equations}{
    name={kinematic equations},
    description={a set of equations that describe the motion of an object under constant acceleration, relating displacement, velocity, acceleration, and time}
}

\newglossaryentry{free fall}{
    name=free fall,
    description={the motion of an object under the influence of gravity alone, with no other forces acting on it}
}





\newglossaryentry{momentum}{
    name={momentum},
    description={the product of a system's mass and velocity}
}

\newglossaryentry{momentum vs. time graph}{
    name={momentum vs.\,time graph},
    description={a graph in which momentum is plotted on the vertical axis and time is plotted on the horizontal axis}
}

\newglossaryentry{impulse}{
    name={impulse},
    description={average net external force multiplied by the time the force acts; equal to the change in momentum}
}

\newglossaryentry{impulse-momentum theorem}{
    name={impulse-momentum theorem},
    description={the impulse equals change in momentum}
}

\newglossaryentry{kinetic energy}{
    name={kinetic energy},
    description={energy of motion}
}

\newglossaryentry{joule}{
    name=joule,
    description={the metric unit for work and energy; equal to 1 newton meter ($\text{N}\cdot\text{m}$)}
}

\newglossaryentry{work}{
    name={work},
    description={force multiplied by distance}
}

%...UNIT 5: FORCE ANALYSIS


\newglossaryentry{contact force}{
    name={contact force},
    description={a type of force that occurs when objects are physically in contact with each other}
}



%...UNIT 7: MOTION IN TWO DIMENSIONS

\newglossaryentry{projectile}{
    name={projectile},
    description={an object that travels through the air and experiences only acceleration due to gravity}
}

\newglossaryentry{projectile motion}{
    name={projectile motion},
    description={the motion of an object that is subject only to the acceleration of gravity}
}

\newglossaryentry{trajectory}{
    name={trajectory},
    description={the path of a projectile through the air}
}

\newglossaryentry{apex}{
    name={apex},
    description={the location on the trajectory at which the projectile reaches maximum height}
}

\newglossaryentry{hang time}{
    name={hang time},
    description={the amount of time that a projectile is in the air during projectile motion}
}

\newglossaryentry{horizontally launched projectile}{
    name={horizontally launched projectile},
    description={a projectile whose initial velocity is entirely in the horizontal direction}
}

\newglossaryentry{impact speed}{
    name={impact speed},
    description={the speed at which a projectile strikes the ground after being launched}
}

\newglossaryentry{circular motion}{
    name={circular motion},
    description={the motion of an object traveling along a circular path at a constant speed}
}

\newglossaryentry{centripetal force}{
    name={centripetal force},
    description={a net force directed toward the center of a circular path that keeps an object moving in a circle}
}

\newglossaryentry{Newton's universal law of gravitation}{
    name={Newton's universal law of gravitation},
    description={states that gravitational force between two objects is directly proportional to the product of their masses and inversely proportional to the square of the distance between them}
}

\newglossaryentry{gravitational constant}{
    name={gravitational constant},
    description={the proportionality constant in Newton's law of universal gravitation}
}

\newglossaryentry{weight}{
    name={weight},
    description={the force of gravity, $F_g$, acting on an object of mass $m$}
}


%...UNIT 8: CONSERVATION IN MECHANICAL SYSTEMS

% \newglossaryentry{system}{
%     name={system},
%     description={one or more objects of interest for which only the forces acting on them from the outside are considered, but not the forces acting between them or inside them}
% }

\newglossaryentry{energy}{
    name={energy},
    description={the ability to do work}
}

\newglossaryentry{potential energy}{
    name={potential energy},
    description={stored energy}
}

\newglossaryentry{gravitational potential energy}{
    name={gravitational potential energy},
    description={energy acquired by doing work against gravity}
}

\newglossaryentry{law of conservation of energy}{
    name={law of conservation of energy},
    description={states that energy is neither created nor destroyed}
}

\newglossaryentry{mechanical energy}{
    name={mechanical energy},
    description={kinetic plus potential energy}
}

\newglossaryentry{elastic collision}{
    name={elastic collision},
    description={a collision in which objects separate after impact and kinetic energy is conserved}
}

\newglossaryentry{inelastic collision}{
    name={inelastic collision},
    description={a collision in which kinetic energy is not conserved}
}

\newglossaryentry{isolated system}{
    name={isolated system},
    description={system in which the net external force is zero}
}

\newglossaryentry{law of conservation of momentum}{
    name={law of conservation of momentum},
    description={when the net external force is zero, the total momentum of the system is conserved or constant}
}

\newglossaryentry{perfectly inelastic collision}{
    name={perfectly inelastic collision},
    description={collision in which objects stick together after impact and kinetic energy is not conserved}
}

\newglossaryentry{recoil}{
    name={recoil},
    description={backward movement of an object caused by the transfer of momentum from another object in a collision}
}

%...UNIT 9: OSCILLATORY MOTION

\newglossaryentry{oscillate}{
    name={oscillate},
    description={to move back and forth regularly between two points}
}

\newglossaryentry{oscillatory motion}{
    name={oscillatory motion},
    description={motion that repeats itself back and forth over the same path, centered around an equilibrium position}
}

\newglossaryentry{equilibrium position}{
    name={equilibrium position},
    description={the location where the object rests if the object is not acted on by a net force}
}

\newglossaryentry{restoring force}{
    name={restoring force},
    description={the force pulling the object back toward equilibrium whenever it moves away}
}

\newglossaryentry{period}{
    name={period},
    description={the time it takes to complete one oscillation cycle}
}

\newglossaryentry{frequency}{
    name={frequency},
    description={the number of oscillation cycles that occur per second}
}

\newglossaryentry{amplitude}{
    name={amplitude},
    description={the maximum distance moved from the equilibrium position}
}

\newglossaryentry{pendulum}{
    name={pendulum},
    description={an object with a small mass suspended from a light wire or string}
}


%...UNIT 12: CONSERVATION OF CHARGE

\newglossaryentry{electric charge}{
    name={electric charge},
    description={a positive or negative property of an object that causes it to be attracted toward or repelled from another charged object; 
    %each charged object generates and is influenced by a force called an electromagnetic force
    }
}

\newglossaryentry{elementary charge}{
    name={elementary charge},
    description={the smallest observed unit of charge that can be isolated in nature; also, the magnitude of charge on 1 proton or 1 electron}
}

\newglossaryentry{electron}{
    name={electron},
    description={subatomic particle that carries one indivisible unit of negative electric charge}
}

\newglossaryentry{proton}{
    name={proton},
    description={subatomic particle that carries the same magnitude charge as the electron, but its charge is positive}
}

\newglossaryentry{electric field}{
    name={electric field},
    description={defines the force per unit charge at all locations in space around a charge distribution}
}

\newglossaryentry{law of conservation of charge}{
    name={law of conservation of charge},
    description={states that total charge is constant in any process}
}

\newglossaryentry{polarization}{
    name={polarization},
    description={separation of charge induced by nearby excess charge}
}

\newglossaryentry{Coulomb's law}{
    name={Coulomb's law},
    description={describes the electrostatic force between charged objects, which is proportional to the charge on each object and inversely proportional to the square of the distance between the objects}
}

\newglossaryentry{electric circuit}{
    name={electric circuit},
    description={physical network of paths through which electric current can flow}
}

\newglossaryentry{circuit component}{
    name={circuit component},
    description={an individual part of an electrical circuit that controls or influences the flow of electric current}
}

\newglossaryentry{simple circuit}{
    name={simple circuit},
    description={a circuit with a single voltage source and a single resistance source through which current flows}
}

\newglossaryentry{electric current}{
    name={electric current},
    description={the flow of electric charge}
}


\newglossaryentry{Ohm's law}{
    name={Ohm's law},
    description={electric current is proportional to the voltage applied across a circuit or other path}
}

\newglossaryentry{resistance}{
    name={resistance},
    description={how much a circuit component opposes the passage of electric current}
}

\newglossaryentry{resistor}{
    name={resistor},
    description={circuit component that provides a known resistance}
}

% \newglossaryentry{potential difference (or voltage)}{
%     name={potential difference (or voltage)},
%     description={change in potential energy of a charge moved from one point to another, divided by the charge; units of potential difference are joules per coulomb, known as volt}
% }

\newglossaryentry{voltage1}{
    name={voltage},
    description={the electrical potential energy per unit charge; electric pressure created by a power source, such as a battery}
}

\newglossaryentry{voltage2}{
    name={voltage},
    description={the electric pressure created by a power source, such as a battery}
}


\newglossaryentry{electric power}{
    name={electric power},
    description={rate at which electric energy is transferred in a circuit}
}

\newglossaryentry{equivalent resistor}{
    name={equivalent resistor},
    description={resistance of a single resistor that is the same as the combined resistance of a group of resistors}
}

\newglossaryentry{in series}{
    name={in series},
    description={when components in a circuit are connected one after the other in the same branch of the circuit}
}

\newglossaryentry{in parallel}{
    name={in parallel},
    description={when a group of resistors are connected side by side, with the top ends of the resistors connected together by a wire and the bottom ends connected together by a different wire}
}

\newglossaryentry{induction}{
    name={induction},
    description={creating an unbalanced charge distribution in an object by moving a charged object toward it (but without touching)}
}

%...UNIT 10: ELECTROMAGNETIC INDUCTION

\newglossaryentry{magnetic dipole}{
    name={magnetic dipole},
    description={term that describes magnets because they always have two poles: north and south}
}

\newglossaryentry{magnetic field}{
    name={magnetic field},
    description={directional lines around a magnetic material that indicates the direction and magnitude of the magnetic force}
}

\newglossaryentry{magnetic pole}{
    name={magnetic pole},
    description={part of a magnet that exerts the strongest force on other magnets or magnetic material}
}

\newglossaryentry{electromagnetic induction}{
    name={electromagnetic induction},
    description={rate at which energy is drawn from a source per unit current flowing through a circuit}
}


\newglossaryentry{Faraday's law}{
    name={Faraday's law},
    description={the means of calculating the emf in a coil due to changing magnetic flux}
}

\newglossaryentry{electromagnet}{
    name={electromagnet},
    description={device that uses electric current to make a magnetic field}
}

\newglossaryentry{transformer}{
    name={transformer},
    description={device that transforms voltages from one value to another}
}

\newglossaryentry{electric motor}{
    name={electric motor},
    description={device that transforms electrical energy into mechanical energy}
}

\newglossaryentry{generator}{
    name={generator},
    description={device that transforms mechanical energy into electrical energy}
}

%...UNIT 11: SIMPLE HARMONIC MOTION & WAVES

\newglossaryentry{wave}{
    name={wave},
    description={a disturbance that moves from its source and carries energy}
}

\newglossaryentry{wave velocity}{
    name={wave velocity},
    description={speed at which the disturbance moves; also called the propagation velocity or propagation speed}
}

\newglossaryentry{wavelength}{
    name={wavelength},
    description={distance between adjacent identical parts of a wave}
}

\newglossaryentry{wave cycle}{
    name={wave cycle},
    description={any portion of a wave encompassed by 1 wavelength}
}

\newglossaryentry{transverse wave}{
    name={transverse wave},
    description={a wave in which the disturbance is perpendicular to the direction of propagation}
}

\newglossaryentry{medium}{
    name={medium},
    description={the solid, liquid, or gas material through which a wave propagates}
}

\newglossaryentry{mechanical wave}{
    name={mechanical wave},
    description={wave that requires a medium through which it can travel}
}

\newglossaryentry{longitudinal wave}{
    name={longitudinal wave},
    description={wave in which the disturbance is parallel to the direction of propagation}
}

\newglossaryentry{sound}{
    name={sound},
    description={a form of energy that travels as a longitudinal mechanical wave through a medium due to vibrating particles}
}

\newglossaryentry{pulse}{
    name={pulse},
    description={a single, short disturbance that travels through a medium}
}

\newglossaryentry{constructive interference}{
    name={constructive interference},
    description={when two waves arrive at the same point exactly in phase; that is, the crests of the two waves are precisely aligned, as are the troughs}
}

\newglossaryentry{destructive interference}{
    name={destructive interference},
    description={when two identical waves arrive at the same point exactly out of phase that is precisely aligned crest to trough}
}






\newglossaryentry{simple harmonic motion}{
    name={simple harmonic motion},
    description={the oscillatory motion in a system where the net force can be described by Hooke’s law}
}

\newglossaryentry{simple harmonic oscillator}{
    name={simple harmonic oscillator},
    description={a device that oscillates in SHM,  such as a mass that is attached to a spring, where the restoring force is proportional to the displacement and acts in the direction opposite to the displacement}
}


\newglossaryentry{electromagnetic wave}{
    name={electromagnetic wave},
    description={a radiant energy wave that consists of oscillating electric and magnetic fields}
}

\newglossaryentry{electromagnetic radiation}{
    name={electromagnetic radiation},
    description={radiant energy that consists of oscillating electric and magnetic fields}
}

